\section{Yor / Marin / Bocchi / Power / Maomao — современные архетипические гибриды}

Последний пункт галереи — не один персонаж, а созвездие недавних архетипических гибридов, которые стоит увидеть вместе, потому что по отдельности каждый лишь один пример, а вместе они показывают тенденцию: современные персонажи всё чаще комбинируют оси, которые раньше считались взаимоисключающими.

Йор Форджер из "Spy x Family" — killer и mother одновременно, agency высочайшего уровня (профессиональный убийца) внутри домашней, почти yamato nadeshiko формы (4.2); гибрид работает, потому что оба слоя поданы всерьёз, не как шутка друг над другом.

Марин Китагава из "My Dress-Up Darling" — genki-энергия (4.1) без хаотичности, направленная в конкретное ремесло (косплей), что делает её редким случаем genki girl с dojikko-партнёром (4.6), а не за счёт собственной неуклюжести.

Бокка из "Bocchi the Rock!" — dandere (3.5), доведённая до предела социальной тревоги, но с ясной творческой agency (музыка), которая существует отдельно от её застенчивости, а не подавлена ею.

Павер из "Chainsaw Man" — хаос (2.7) без tsundere-смягчения, редкий случай, где genki-подобная энергия соединена с буквально нечеловеческой моралью, а не с обычным добрым сердцем под колючей внешностью.

Маомао из "Kusuriya no Hitorigoto" ("The Apothecary Diaries") — meganekko-интеллект (4.5) без застенчивости, соединённый с почти патологической, но продуктивной жадностью до головоломок, которая работает как её собственная версия контрастного бага (глава 4): гениальный аналитик, которая ставит любопытство выше собственной безопасности.

Общий вывод для книги: ни один из этих персонажей не собран из единственного архетипа — все они гибриды, и именно гибридность делает их современными, а не устаревшими версиями более чистых типов из глав 3--4. Для RIN это подтверждение общего курса книги: её собственный вектор (глава 12) должен быть таким же осознанным гибридом, а не выбором одного архетипа из каталога.
